\numberlesssection{ВВЕДЕНИЕ}

Современные системы нередко строятся на основе микросервисной~\cite{ref:microservice-arch} 
архитектуры, состоящей из различных, связанных друг с другом, компонентов, например
микросервисов, баз данных, внешних API, брокеров сообщений и других.
В масштабных системах количество связей между компонентами может быть огромным.

Одним из ключевых требований к любой программной системе является надежность
и отказоустойчивость. При большом количестве связей ручной анализ зависимостей
компонентов друг от друга становится достаточно кропотливой задачей.
Решением данной проблемы является автоматизация поиска зависимостей конкретных
компонентов и симуляция каскадных сбоев, что сведет к минимуму вероятность ошибок
вследствие человеческого фактора и уменьшит временные затраты на проектирование
и анализ архитектуры системы.

Цель курсовой работы --- разработать базу данных для анализа отказоустойчивости
систем на основе микросервисной архитектуры.
Для достижения поставленной цели необходимо решить следующие задачи:
\begin{itemize}
    \item провести анализ предметной области;
    \item провести анализ существующих решений;
    \item сформулировать требования и ограничения к разрабатываемой базе данных;
    \item спроектировать ролевую модель на уровне приложения;
    \item выбрать СУБД для реализации разрабатываемой базы данных;
    \item спроектировать сущности базы данных и ограничения целостности топологии сети;
    \item сформулировать требования и ограничения к разрабатываемому приложению;
    \item выбрать средства реализации базы данных и приложения;
    \item описать методы тестирования разработанного функционала и разработать тесты
        для проверки корректности работы приложения;
    \item исследовать зависимость времени выполнения запросов поиска зависимостей
        от количества узлов системы.
\end{itemize}
